\documentclass[10pt,a4paper]{amsart}
\usepackage[margin=19mm]{geometry}
\usepackage{amsmath,amssymb,mathtools}
\usepackage[scr=rsfs,cal=euler]{mathalfa}
\usepackage{xfrac}
\usepackage[unicode,hidelinks,bookmarks=false]{hyperref}
\newcommand{\ie}{i.e.\ }
\newcommand{\cf}{cf.\ }
\newcommand{\resp}{respectively\ }
\newcommand{\Subsection}{\subsection}
\providecommand{\nobreakdash}{\nobreak\hbox{-}}
\setlength{\emergencystretch}{2em}
\multlinegap0pt
\usepackage{bm}
\usepackage{bbm}
\usepackage{dsfont}
\usepackage{xspace}
\usepackage{cancel}
\usepackage{mathtools}
\usepackage{amsthm}
\makeatletter
\@ifundefined{thm}{\newtheorem{thm}{Thm}}{}
\@ifundefined{prop}{\newtheorem{prop}[thm]{Proposition}}{}
\makeatother
\DeclareMathOperator{\tpmd}{tpmd}
\title[(Almost) Complete Intersection Lov\'{a}sz-Saks-Schrijver ide]{(Almost) Complete Intersection Lov\'{a}sz-Saks-Schrijver ideals and regularity of their powers}
\author{Marie Amalore Nambi, Neeraj Kumar, Chitra Venugopal}
\date{Source: arXiv:2305.09190v4 (CC BY 4.0)}
\begin{document}
\maketitle

\noindent\emph{Source and licence.} (Almost) Complete Intersection Lov\'{a}sz-Saks-Schrijver ideals and regularity of their powers. Version arXiv:2305.09190v4 (CC BY 4.0), distributed under the Creative Commons Attribution 4.0 International licence, as recorded in the arXiv licence metadata for this exact version and shown on the abstract page https://arxiv.org/abs/2305.09190v4. Full rights notice: licenses/arxiv-2305.09190-CC-BY-4.0.txt.\par\smallskip

\noindent\textit{Editorial scope.} Counted unit: the Prop whose source label is \texttt{LT} in the article (source0.tex lines 620--628, source label \texttt{LT}), delivered together with the article's own complete original proof of it (source0.tex lines 629--666, one \texttt{proof} environment of 38 lines). No mathematics is written, repaired or re-derived: statement and proof are the article's own text line for line. Required context retained uncounted: none. Every definition, display and label of those lines is retained unchanged. Omitted from this package and not redistributed: 1--619, 667--1278. Nothing inside a retained excerpt is omitted or altered: every retained body is delivered exactly as the article's lines are written, so no omission receipt of any kind accompanies this package. Citations: the retained excerpts carry no citation command at all, so no bibliography entry of the article is transcribed and no reference is redistributed. Reference adaptation: none was needed and none was made; every reference inside the delivered unit resolves inside it, so no unresolved reference or citation is produced. Printed slips kept literal: no printed word, symbol, hypothesis or number of the counted statement or of its proof was altered. Layout and class: the article's own class is \texttt{amsart} with packages including inputenc, amsmath, geometry, amsfonts, amssymb, amsthm, dsfont, amsxtra, enumerate, graphicx, xcolor, hyperref, multicol, mathtools; the wrapper is the lane's pinned \texttt{amsart} editorial wrapper at 10pt on A4, which restates the theorem-like environments the retained text needs and adds the article's own macro definitions that the retained text needs (\texttt{\textbackslash tpmd}), with extra wrapper packages (bm, bbm, dsfont, xspace, cancel, mathtools). The delivered numbering is the unit's own. Assets: no retained span contains a figure environment, a graphics-inclusion command, a PDF-page inclusion, a table or a TikZ picture; no image file is copied into this package, the package declares no asset dependency and no image file is redistributed with it. Licence: version arXiv:2305.09190v4 of this article is distributed under the Creative Commons Attribution 4.0 International licence, as recorded in the arXiv licence metadata for this exact version and shown on the abstract page https://arxiv.org/abs/2305.09190v4. A full rights notice is delivered at licenses/arxiv-2305.09190-CC-BY-4.0.txt. Characters: every retained excerpt is pure ASCII, so the delivered engine, XeLaTeX with its default Latin Modern text font, reports no missing character.
\begin{prop} \label{LT}
     Let $G=(V(G),E(G))$ be a graph, $d \geq p=\tpmd(G)$ and $E(G)= \cup_{\ell=1}^{2p}M_{\ell}$ be a twisted matching decomposition. If $G$ admits a twisted positive matching decomposition with respect to twisted matching decomposition $E(G)=\cup_{\ell=1}^{2p}M_\ell$, then there exists a term order $<$ on $S$ such that for all  $q=1,2,\ldots,p$, for every $A=\{ i,j \} \in (M_{2q-1},M_{2q})$,
    \begin{equation} \label{LT-eq}
     \text{in}_<(\hat{f}_A^{(d)})= \begin{cases}
         x_{i2q-1}x_{j2q} \hspace{5mm} \text{ if } \{i,j\} \in M_{2q-1}\\
         x_{i2q}x_{j2q-1} \hspace{5mm} \text{ if } \{ i,j \} \in M_{2q}.
     \end{cases}
    \end{equation}
\end{prop}

\begin{proof}
Consider $E(G)= \cup_{\ell=1}^{2p}M_{\ell}$ be the matching decomposition of $G$ and $E(G) = \cup_{q=1}^p(M_{2q-1},M_{2q})$ be a twisted matching decomposition with the respective weight functions $w_{q}:V(H_q) \rightarrow \mathbb{R}$. In order to define the required term order $<$, we define the weight vectors $\mathfrak{w}_1, \ldots, \mathfrak{w}_p \in \mathbb{R}^{|V| \times 2d}$ as follows,

\begin{itemize}
    \item $\mathfrak{w}_{q}(x_{ik})=0$ if $k \notin \{2q-1,2q\}$ and
    \item $\mathfrak{w}_{q}(x_{ik})=w_q(i_k)$ if $k \in \{2q-1,2q\}$.
\end{itemize}

Then, by construction, it follows that:
\begin{equation} \label{intial1}
    \text{in}_{\mathfrak{w}_{1}}(\hat{f}_{\{i,j\}}^{(d)})= \begin{cases}
    x_{i1}x_{j2} \hspace{5mm} \text{ if } \{i,j\} \in M_{1}, \\
    x_{i2}x_{j1} \hspace{5mm} \text{ if } \{i,j\} \in M_{2}, \\
    \sum_{k=2}^d(x_{i2k-1}x_{j2k}-x_{i2k}x_{j2k-1}) \hspace{5mm} \text{ if } \{i,j\} \notin (M_{1},M_2).
\end{cases}
\end{equation} 


With the weight vectors defined, we say $x^{\alpha} < x^{\beta}$ if,
\begin{enumerate}
    \item $|\alpha|<|\beta|$ or 
    \item $|\alpha|=|\beta|$ and $\mathfrak{w}_{q}(x^{\alpha})<\mathfrak{w}_{q}(x^{\beta})$ for the smallest $q$ such that $\mathfrak{w}_{q}(x^{\alpha})\neq \mathfrak{w}_{q}(x^{\beta}) $ or
    \item $|\alpha|=|\beta|$ and $\mathfrak{w}_{q}(x^{\alpha})=\mathfrak{w}_{q}(x^{\beta})$ for all $q$ and $x^{\alpha}<_0x^{\beta}$ for an arbitrary but fixed term order $<_0$.
\end{enumerate}

For a given edge $A=\{i,j\} \in E(G)$, one has $A \in (M_{2q-1},M_{2q})$ for some $1\leq q \leq p$. There are two possibilities, either $A \in (M_{1},M_{2})$ or $A \notin (M_{1},M_{2})$.

If $A \in (M_{1},M_{2})$, then the statement follows from Equation (\ref{intial1}).

If $A \notin (M_{1},M_{2})$, then, from the defined weight vector $\mathfrak{w}_{1}$, it follows that $\mathfrak{w}_{1}(x_{i1}x_{j2})<\mathfrak{w}_{1}(y)$ and $\mathfrak{w}_{1}(x_{i2}x_{j1})<\mathfrak{w}_{1}(y)$ for all $y \in \{x_{i,2k-1}x_{j,2k},x_{i,2k}x_{j,2k-1}\}$, where $k=2,\ldots, d$. This implies that $x_{i1}x_{j2}$ and $x_{i2}x_{j1}$ are not the leading term of $\hat{f}_A^{(d)}$ with respect to $<$, and so the leading term is a monomial in $\sum_{k=2}^d(x_{i2k-1}x_{j2k}-x_{i2k}x_{j2k-1})$. Now, in the next step, there are two possibilities, either $A \in (M_{3},M_{4})$ or $A \notin (M_{3},M_{4})$. 

If $A \in (M_{3},M_{4})$, then, from the weight vector $\mathfrak{w}_{2}$, it follows that $x_{i3}x_{j4} (x_{i4}x_{j3})$ is the leading term of $\hat{f}_A^{(d)}$ if $A \in M_3 (M_4)$, respectively. 

If $A \notin (M_{3},M_{4})$, then, from the weight vector $\mathfrak{w}_{2}$, it follows that $\mathfrak{w}_{2}(x_{i3}x_{j4})<\mathfrak{w}_{2}(y)$ and $\mathfrak{w}_{2}(x_{i4}x_{j3})<\mathfrak{w}_{2}(y)$ for all $y \in \{x_{i,2k-1}x_{j,2k},x_{i,2k}x_{j,2k-1}\}$, where $k=3,\ldots, d$. This implies that with respect to the monomial order $<$, $x_{i3}x_{j4}$ and $x_{i4}x_{j3}$ are not the leading term of $\hat{f}_A^{(d)}$, and so the leading term is a monomial in $\sum_{k=3}^d(x_{i2k-1}x_{j2k}-x_{i2k}x_{j2k-1})$. 
In this way, by repeating the above process one obtains the desired result. 

Observe that this process terminates in a finite number of steps as the leading term of $\hat{f}_A^{(d)}$ will be obtained at the $q$th step, where $q \leq p$ is the integer such that $A \in (M_{2q-1},M_{2q})$.
\end{proof}

\end{document}
